\newpage

\rhead{\bfseries \thepage}
\vspace*{\fill}
\begin{center}
{\bf Abstract}
\end{center}
\addcontentsline{toc}{section}{Abstract}
This thesis gives an introduction to {\tt Quantum Computation} and discusses implementation strategies for a {\tt Quantum Computer} simulator on a distributed system. Research into {\tt Quantum Computers} has expanded rapidly since Peter Shor showed that integer factorisation --- believed to be classically intractable --- admits a polynomial-time quantum algorithm; further quantum algorithms have since been discovered that outperform their classical counterparts on a range of problems. Physical realisation of {\tt Quantum Computers} remains difficult, and machines built to date are limited in size and noisy in operation; as a result, much of the research into {\tt Quantum Computation} algorithms makes use of simulators running on classical hardware. The fundamental difference between classical and quantum computation makes such simulation inherently expensive: the state vector grows exponentially in the number of qubits. This thesis examines the strategies available for simulating a {\tt Quantum Computer} on a parallel {\tt Linux} cluster, identifies the central scaling pitfall of dense-matrix approaches, and presents the in-place sparse-gate strategy that scalable simulators must adopt instead. It then applies that strategy to the three canonical quantum algorithms: the {\tt Quantum Fourier Transform}, Grover's unstructured search, and Shor's integer factorisation.
\vspace*{\fill}
